\chapter{总结与展望}

\section{工作总结}

本文围绕基于时空特征融合和稀疏表示的 GNSS NLOS 识别方法展开研究，主要完成以下工作：

（1）提出 STCA 模型，设计时空双通道特征提取架构。针对 GNSS 信号数据同时包含空间分布特性和时序演化规律的特点，设计了并行的时空特征提取路径，分别捕捉卫星空间几何关系与信号质量随时间的变化模式。

（2）实现交叉注意力融合模块与稀疏表示层。通过交叉注意力机制实现时空特征的双向信息交互与自适应融合，并利用稀疏表示层对融合后的高维特征进行紧凑编码，在降低冗余的同时增强特征的判别性。

（3）在内域场景下取得 73.89\% 的识别准确率，训练损失收敛至 0.6273。实验结果表明，STCA 模型能够有效区分视距信号与非视距信号，性能优于传统机器学习方法。

（4）通过完整的消融实验验证各模块的贡献。实验分别考察了空间编码器、时序编码器、交叉注意力模块及稀疏表示层对模型性能的独立影响，为架构设计的合理性提供了实证支持。时间窗口大小敏感性分析表明，窗口大小为 10 时模型性能最优。

\section{存在不足与未来方向}

本研究仍存在以下局限，后续工作可从以下方面继续推进：

（1）外域泛化性能有待提升。当前模型在跨地点场景下的表现较内域场景有所下降，提升模型对未见环境的适应能力是实际应用面临的主要挑战。未来可探索域对抗训练、元学习等域泛化方法。

（2）时序模块的建模能力尚未充分挖掘。当前采用的 LSTM 模块虽能捕捉基本的时序依赖关系，但对长程时间动态的建模能力有限。Transformer 架构、时间卷积网络等更强大的时序模型值得进一步研究。

（3）多模态数据融合与在线学习策略具有研究价值。实际 GNSS 接收环境复杂多变，除信号特征外，还可融合视觉、惯性传感器等多源信息。同时，研究增量学习与在线更新机制，将使模型能够适应环境变化，降低维护成本。

\section{主要创新点}

本文的主要创新点归纳如下：

（1）提出了适用于 GNSS NLOS 识别的时空双通道架构。现有深度学习模型要么仅处理单历元空间特征，要么仅处理时序序列，未能同时利用 GNSS 信号的空间分布和时间演化特性。STCA 模型通过并行的空间编码器和时序编码器，实现了对两类信息的均衡建模。

（2）设计了交叉注意力时空融合机制。不同于简单的特征拼接或加权平均，交叉注意力机制使模型能够根据时序状态自适应地关注相关的空间特征，建立了空间几何分布与信号时间演变之间的动态关联。

（3）引入稀疏表示层增强特征判别性。在分类器前增加稀疏性约束，迫使模型学习更紧凑的特征表示，在降低冗余的同时提升了对噪声和异常值的鲁棒性。

（4）完成了系统的实验验证。在内域和外域两种场景下评估模型性能，通过对比实验、消融实验、统计显著性检验等多种方式，全面验证了 STCA 模型的有效性和优越性。

\section{研究局限}

客观分析，本研究存在以下局限性：

（1）数据规模有限。实验数据来自 7 个测点的静态采集，总样本量约 XXXX 个。虽然数据质量较高，但地域覆盖范围有限，可能影响模型的泛化能力。

（2）场景单一。实验数据均为静态采集，未涉及车载、行人等动态场景。动态环境中的多普勒效应、接收机运动等因素可能对信号特征产生影响，当前模型未经过此类场景的验证。

（3）实时部署验证不足。虽然模型复杂度分析表明 STCA 属于轻量级模型，推理延迟小于 10ms，但尚未在实际接收机或移动设备上进行部署验证。边缘设备的内存限制、功耗约束等因素可能影响实际性能。

（4）可解释性仍有提升空间。尽管引入稀疏表示层增强了特征判别性，但深度学习模型的"黑箱"本质未变。注意力权重的可视化分析仅提供有限的事后解释，无法完全揭示模型的决策机制。

\section{未来研究方向}

基于当前研究的局限性，未来工作可从以下方向继续推进：

（1）数据扩充与增强。收集更多地测点的数据，涵盖不同城市、不同时间段、不同天气条件。探索数据增强技术，如添加噪声、时间扭曲、特征混合等，在不增加采集成本的前提下扩充训练数据。

（2）动态场景验证。在车载、行人等动态场景中采集数据，验证 STCA 模型的适用性。针对动态场景的特点，可考虑引入速度、加速度等运动状态特征。

（3）域泛化方法研究。探索域 adversarial 训练、元学习、测试时自适应等域泛化技术，提升模型对未见环境的适应能力。目标是在新地点无需重新训练即可获得较好的性能。

（4）模型压缩与加速。研究知识蒸馏、剪枝、量化等模型压缩技术，在保持性能的前提下进一步降低模型复杂度和推理延迟，满足边缘设备的部署需求。

（5）多星座融合。当前研究仅使用 GPS 单星座数据。未来可扩展到多星座融合（GPS、BDS、Galileo、GLONASS），利用更多卫星资源提升 NLOS 识别性能。

（6）可解释性增强。探索注意力可视化、梯度分析、特征重要性评估等可解释性方法，深入理解模型的决策机制，为模型改进提供指导。

\section{结语}

GNSS NLOS 信号识别是提升城市环境定位精度的关键技术。本文提出的 STCA 模型在时空特征融合和稀疏表示方面进行了有益探索，取得了一定的研究成果。但由于研究时间和个人能力限制，工作中难免存在不足之处。希望本研究能为 GNSS NLOS 检测领域的发展贡献一份力量，也期待后续研究者在此基础上取得更大突破。

\section{本章小结}

本章总结了全文的研究工作与创新点，客观分析了当前方法的局限性，包括数据规模有限、场景单一、实时部署验证不足、可解释性有待提升等。未来研究方向包括数据扩充与增强、动态场景验证、域泛化方法研究、模型压缩与加速、多星座融合、可解释性增强等。
